\documentclass[11pt]{article}

\usepackage[a4paper,margin=1in]{geometry}
\usepackage[T1]{fontenc}
\usepackage[utf8]{inputenc}
\usepackage{lmodern}
\usepackage{graphicx}
\usepackage{hyperref}

\hypersetup{
  colorlinks=true,
  linkcolor=black,
  urlcolor=blue,
  pdftitle={Stability Oracle Documentation Paper},
  pdfauthor={OpenAI Codex},
  pdfsubject={Structural Stability Evaluation Layer for HAOS-IIP}
}
\setcounter{secnumdepth}{3}
\setcounter{tocdepth}{2}

\title{Stability Oracle\\Documentation Paper}
\author{Structural Stability Evaluation Layer for HAOS-IIP}
\date{March 24, 2026}

\begin{document}
\maketitle
\tableofcontents
\newpage

\section{Purpose}

Stability Oracle is a minimal structural evaluation engine for HAOS-IIP-aligned systems.
Its role is narrow by design:

\begin{itemize}
\item evaluate whether structure survives controlled perturbation
\item expose deterministic stability signals
\item provide a machine-readable decision surface for tools and agents
\end{itemize}

The oracle does not model physics, derive ontology, or make cosmological claims.
It acts as an operational measurement layer over a frozen external HAOS telemetry surface.

\section{System Framing}

\subsection{What The Oracle Measures}

The oracle evaluates structural survival across four bounded metrics:

\begin{itemize}
\item \texttt{structural\_retention}: how much organization persists after perturbation
\item \texttt{temporal\_consistency}: whether ordering relations remain recoverable
\item \texttt{causal\_deformation}: how strongly propagation geometry is distorted
\item \texttt{geometric\_integrity}: whether local coherence collapses or stabilizes
\end{itemize}

These metrics are normalized to the interval \texttt{[0,1]}.
The intent is to support synthetic graphs, interaction networks, symbolic state spaces, agent memory structures, and simulation outputs without coupling the oracle to one substrate.

\subsection{What The Oracle Does Not Claim}

The oracle does not claim:

\begin{itemize}
\item emergence of spacetime
\item physical law derivation
\item consciousness modeling
\item universal ontology
\item predictive cosmology
\end{itemize}

It is a measurement instrument, not a theory.
Its labels such as \texttt{stable}, \texttt{marginal}, and \texttt{unstable} describe structural survivability only.

\section{Architecture}

The standalone repository is organized as a layered system:

\begin{enumerate}
\item agent layer
\item tool interface layer
\item skill logic layer
\item adapter layer
\item external HAOS dependency
\end{enumerate}

No scientific code lives in this repository.
The project treats HAOS-IIP as an external dependency and communicates with it through a bounded CLI adapter.

\begin{figure}[h]
  \centering
  \includegraphics[width=\textwidth]{../../Images/Stability Oracle Evaluation Layer Overview.png}
  \caption{Stability Oracle evaluation layer overview.}
\end{figure}

\section{State Language}

\subsection{Invariant Core}

The invariant contract is defined in \texttt{haos\_skill/state\_spec.py}.
It contains four typed dataclasses:

\begin{itemize}
\item \texttt{State}
\item \texttt{Perturbation}
\item \texttt{StabilityMetrics}
\item \texttt{OracleResult}
\end{itemize}

This layer is intentionally algorithm-free.
It defines the language the oracle speaks.

\subsection{State}

\texttt{State} is a minimal structural representation with:

\begin{itemize}
\item node identifiers
\item directed edges
\item optional feature map
\item optional timestamp map
\end{itemize}

The contract enforces:

\begin{itemize}
\item unique node identifiers
\item edges that reference only known nodes
\item optional feature and timestamp keys that reference only known nodes
\end{itemize}

\subsection{Perturbation}

\texttt{Perturbation} describes bounded structural disturbance:

\begin{itemize}
\item node drop ratio
\item edge drop ratio
\item noise ratio
\item cluster split flag
\end{itemize}

This is descriptive only.
It does not imply any specific perturbation engine implementation.

\subsection{Oracle Result}

\texttt{OracleResult} now includes:

\begin{itemize}
\item classification
\item metrics
\item confidence
\item coherence score
\item raw metric vector
\item normalized metric vector
\item policy version
\end{itemize}

This preserves reproducibility across policy revisions.

\section{Metric Interface}

\subsection{Design Goal}

Metrics are designed to be:

\begin{itemize}
\item deterministic
\item side-effect free
\item reproducible
\item bounded to \texttt{[0,1]}
\item defined only on \texttt{(State\_before, State\_after)}
\end{itemize}

\subsection{Protocol}

The interface is defined in \texttt{haos\_skill/metrics/base.py}.

\begin{verbatim}
class StabilityMetric(ABC):
    name: str
    def compute(self, before: State, after: State) -> float:
        ...
\end{verbatim}

The companion \texttt{MetricRegistry} holds pluggable metrics and computes them as a deterministic bundle.

\subsection{Current Minimal Metrics}

The first concrete metrics are deliberately small:

\begin{itemize}
\item \texttt{StructuralRetentionMetric}
\item \texttt{TemporalConsistencyMetric}
\item \texttt{CausalDeformationMetric}
\item \texttt{GeometricIntegrityMetric}
\end{itemize}

They are meant to be composable, testable, and replaceable.
They are not presented as final truth functions.

\section{Classification Policy}

\subsection{Policy Intent}

Classification is implemented as a deterministic policy module in \texttt{haos\_skill/oracle/classifier.py}.
It is not a learned model.
It is a small, inspectable decision rule.

\subsection{Normalization}

The policy converts the raw metric bundle into a direct stability vector:

\[
\begin{array}{rcl}
S_1 & = & \mbox{structural\_retention} \\
S_2 & = & \mbox{temporal\_consistency} \\
S_3 & = & 1 - \mbox{causal\_deformation} \\
S_4 & = & \mbox{geometric\_integrity}
\end{array}
\]

All normalized dimensions therefore share the same polarity:
higher means more stable.

\subsection{Decision Rule}

Version \texttt{v1\_floor\_mean\_band} applies:

\begin{itemize}
\item a hard survivability floor at \texttt{0.30}
\item an unweighted coherence score equal to the arithmetic mean of the normalized vector
\item a spread score equal to \texttt{max(S) - min(S)}
\end{itemize}

Decision regions:

\begin{itemize}
\item \texttt{stable}: coherence \texttt{>= 0.75} and spread \texttt{<= 0.35}
\item \texttt{marginal}: coherence \texttt{>= 0.50} and no floor violation
\item \texttt{unstable}: every remaining case
\end{itemize}

\subsection{Confidence}

The policy emits a continuous confidence value:

\[
\mbox{confidence} = \mbox{coherence} \times (1 - \mbox{spread})
\]

This gives agents both a discrete verdict and a bounded gradient.

\section{Adapter And Tool Surface}

\subsection{External Dependency Boundary}

The adapter layer calls an external HAOS oracle CLI.
By default it expects:

\begin{verbatim}
haos_iip.demo stability --json
\end{verbatim}

Input is passed either:

\begin{itemize}
\item by substituting a \texttt{\{state\_json\}} token
\item by appending the value after \texttt{--state-json}
\item or through standard input
\end{itemize}

\subsection{Skill Surface}

The public callable interface exposes:

\begin{itemize}
\item \texttt{evaluate\_structure(state\_spec)}
\item \texttt{scan\_structure(grid\_spec)}
\end{itemize}

The CLI exposes:

\begin{verbatim}
haos-skill evaluate state.json
haos-skill scan grid.json
\end{verbatim}

Outputs are JSON only.

\subsection{Agent Integration}

The repository includes minimal integration patterns for:

\begin{itemize}
\item OpenClaw-style skill registration
\item Hermes-style schema plus callable registration
\end{itemize}

\section{Safety And Reproducibility}

The safety layer enforces:

\begin{itemize}
\item hard timeout guards
\item serialized input size limits
\item node, edge, and batch size bounds
\item fail-closed protocol handling
\item optional deterministic in-memory caching
\end{itemize}

The classifier additionally preserves reproducibility with an explicit policy version string.

\section{Current Documentation Inventory}

This paper consolidates the current documentation surface:

\begin{tabular}{p{0.33\textwidth} p{0.57\textwidth}}
\textbf{Document} & \textbf{Role} \\
\hline
\texttt{README.md} & Repository landing page, usage, architecture, and public entry points. \\
\texttt{docs/WHAT\_IS\_STABILITY\_ORACLE.md} & Conceptual framing, scope, non-claims, and roadmap. \\
\texttt{docs/01\_STABILITY\_ORACLE\_DOCUMENTATION\_PAPER.tex} & Numbered documentation source for this compiled paper. \\
\texttt{Images/Stability Oracle Evaluation Layer Overview.png} & Visual architecture overview used on the public landing page. \\
\texttt{Images/Stability\_Oracle.pdf} & Existing companion PDF asset. \\
\end{tabular}

\section{Verification Status}

At the time of writing, the standalone repository includes automated tests for:

\begin{itemize}
\item state contract normalization
\item metric interface behavior
\item classification policy
\item timeout handling
\item deterministic cache behavior
\item skill output shape
\end{itemize}

This confirms that the project currently has:

\begin{itemize}
\item a stable language layer
\item a pluggable metric interface
\item a reproducible classification policy
\item a bounded agent-facing skill surface
\end{itemize}

\section{Roadmap}

The current development sequence is:

\begin{enumerate}
\item language layer
\item metric interface
\item oracle engine
\item public demo surface
\item agent skill integration
\item research extensions
\end{enumerate}

The next clean architectural step is the \texttt{OracleEngine}, which should:

\begin{enumerate}
\item run perturbation
\item run the metric registry
\item map the resulting metrics through the classifier policy
\end{enumerate}

\section{Conclusion}

Stability Oracle advances through small operational steps.
Its value does not depend on large claims.

The system becomes useful by remaining:

\begin{itemize}
\item deterministic
\item bounded
\item testable
\item composable
\item explicit about what it does and does not claim
\end{itemize}

That is the foundation that allows the project to function as a practical bridge between frozen HAOS-IIP research artifacts and usable engineering interfaces.

\end{document}
